\CourseSection{Introduction}\label{sec:introduction}
\Guidance{The body of the completed proposal, \emph{i.e.} Sections 1--6 below, has a maximum length of eight pages; the title page, executive summary, contents, references, and appendices are not included in this limit. \\
In the intro below, establish the wider problem you are trying to address, project's scope in the context of this problem, the general design, and limitations. \\
Change the titles of these subsections as you see fit. We have just provided some example placeholders.}

\CourseSubsection{Inefficiencies in unilateral phase detractors}

\Guidance{Describe the problem, the intended end user or market, evidence of the need, and the value of addressing it. Support the context with several citations from peer-reviewed literature \cite{kraft97rockwellEncabulator}.\footnote{Not exactly a peer-reviewed source, but I wanted to put an example reference somewhere :)} Feel free to change the section titles to fit your project. This section is the 'Problem and Purpose' section.}

\CourseSubsection{Retro-encabulator design for inverse reactive current}\label{sec:project-objectives}
\Guidance{State the specific objectives and measurable outcomes that address the problem. Characterize the electrical and / or computer engineering involved. Support your claims by describing similar approaches in the peer-reviewed literature. Include a general diagram showing the system context. Compare your approach to those in the peer-reviewed literature. This section is the 'Project Objective' section.}

\CourseSubsection{Deliverables: Power production, elimination of side-fumbling, and reduction of sinusoidal depleneration}
\Guidance{This is a refinement of Sec. \ref{sec:project-objectives}. Define the end product or products, the operating context, system inputs and outputs, and the major components. Explain each component's purpose and how the components fit together. State the project boundaries and assumptions. Address several limitations of your approach. Provide citations in the peer-reviewed literature and / or textbooks that apply similar strategies. This section is the 'Scope and deliverables' section.}
\clearpage
